\documentclass[a4paper]{article}

% Español (México) con XeLaTeX: fontspec para fuentes Unicode, polyglossia
% para la división silábica y los nombres del documento en la variante mexicana.
\usepackage{fontspec}
\usepackage{polyglossia}
\setdefaultlanguage[variant=mexican]{spanish}
% Los nombres chinos van como «pinyin (original)»: xeCJK da a los caracteres
% Han su propia fuente; sin ella XeLaTeX los omite con un aviso.
\usepackage{xeCJK}
\setCJKmainfont{FandolSong-Regular.otf}
% polyglossia no traduce los nombres de listings.
\AtBeginDocument{\providecommand{\lstlistingname}{}\renewcommand{\lstlistingname}{Listado}%
  \providecommand{\lstlistlistingname}{}\renewcommand{\lstlistlistingname}{Índice de listados}}
\usepackage{amsmath, amssymb}
\usepackage{graphicx}
\usepackage[margin=2.5cm]{geometry}
\usepackage[most]{tcolorbox}
\usepackage{etoolbox}
\usepackage{listings}
\usepackage{hyperref}
\usepackage{booktabs}
\usepackage{float}

\newtcolorbox{knowledgebox}[1]{enhanced, colback=blue!5!white, colframe=blue!75!black, colbacktitle=blue!75!black, coltitle=white, fonttitle=\bfseries, title={#1}, attach boxed title to top left={yshift=-2mm, xshift=2mm}, boxrule=1pt, sharp corners}
\newtcolorbox{importantbox}[1]{enhanced, colback=yellow!10!white, colframe=yellow!80!black, colbacktitle=yellow!80!black, coltitle=black, fonttitle=\bfseries, title={#1}, sharp corners}
\newtcolorbox{warningbox}[1]{enhanced, colback=red!5!white, colframe=red!75!black, colbacktitle=red!75!black, coltitle=white, fonttitle=\bfseries, title={#1}, sharp corners}

\lstset{
    language=Python,
    basicstyle=\ttfamily\small,
    keywordstyle=\color{blue},
    stringstyle=\color{red!60!black},
    commentstyle=\color{green!60!black},
    breaklines=true,
    frame=single,
    numbers=left,
    numberstyle=\tiny\color{gray},
    captionpos=b,
    extendedchars=true
}

\newcommand{\notetitle}{Plan Mode y update\_plan en Codex}
\newcommand{\noteauthors}{Elaboradas a partir de la serie Modern Agent de Wudaokou Nash (五道口纳什)}
\newcommand{\notedate}{2026}
\newcommand{\videochannel}{Wudaokou Nash (五道口纳什)}
\newcommand{\videopublishdate}{2026-04-09}
\newcommand{\videoduration}{14:49}
\newcommand{\videourl}{https://www.bilibili.com/video/BV1NdDtBjEg7/}
\newcommand{\videocoverpath}{cover.jpg}
\newcommand{\repourl}{https://github.com/hqhq1025/ai-course-notes}

\usepackage{fancyhdr}
\pagestyle{fancy}
\fancyhf{}
\fancyfoot[C]{\thepage}
\fancyfoot[R]{\footnotesize\color{gray}\href{\repourl}{AI Course Notes}}
\renewcommand{\headrulewidth}{0pt}
\renewcommand{\footrulewidth}{0pt}

\begin{document}

\begin{titlepage}
\centering
{\Large Notas del curso\par}\vspace{1.2cm}
{\huge\bfseries \notetitle\par}\vspace{0.8cm}
{\large \noteauthors\par}\vspace{0.3cm}
{\large \notedate\par}\vspace{1.2cm}
\ifdefempty{\videocoverpath}{}{\includegraphics[width=0.82\textwidth,height=0.45\textheight,keepaspectratio]{\videocoverpath}\par}
\vfill
\begin{tcolorbox}[width=0.9\textwidth, colback=black!2!white, colframe=black!60, sharp corners]
\textbf{Autor o canal del video}: \videochannel\par
\textbf{Fecha de publicación}: \videopublishdate\par
\textbf{Duración del video}: \videoduration\par
\ifdefempty{\videourl}{}{\textbf{Enlace del video}: \href{\videourl}{\nolinkurl{\videourl}}\par}
\textbf{Página del proyecto}: \href{\repourl}{\nolinkurl{\repourl}}\par
\end{tcolorbox}
\end{titlepage}

\tableofcontents
\newpage

\section{Introducción: por qué la planeación es una capacidad central del coding agent}

Esta clase continúa la discusión sobre el diseño de herramientas para coding agents del tipo Codex / Claude Code. Quien da la clase señala que, al presentar antes el conjunto mínimo de herramientas de Codex, es fácil pasar por alto una herramienta que parece simple pero es muy importante: \texttt{update\_plan}. En una tarea de desarrollo real, un agent no enfrenta una sola pregunta con una sola respuesta, sino una tarea sobre un repositorio de código que se despliega de forma continua: primero entender el codebase, después formar un plan, luego ejecutar, verificar y ajustar, hasta terminar la tarea.

\begin{figure}[H]
\centering
\includegraphics[width=\textwidth]{frames/frame-00-00-50.jpg}
\caption{En la serie de Modern Agent, el Planning Agent y el diseño de herramientas de Codex forman un hilo continuo.\protect\footnotemark}
\end{figure}
\footnotetext{Intervalo de tiempo en el video: 00:00:45--00:00:55.}

\begin{importantbox}{Flujo típico de tareas de Codex}
En la abstracción de quien da la clase, Codex procesa una necesidad de desarrollo normalmente en tres pasos: \textbf{codebase exploration}, \textbf{plan} y \textbf{execution}. La exploration se encarga de leer el código y el contexto de forma intensiva, el plan divide el objetivo en pasos ejecutables, y la execution avanza de manera continua según la retroalimentación de la ejecución.
\end{importantbox}

Esta es también la razón por la que la planeación no es solo una función de la interfaz. En realidad es una representación de estado dentro del agent loop: el modelo necesita conservar en el contexto «qué hay que hacer ahora», «hasta qué paso se ha llegado» y «si hace falta replanear», y mostrarlo al usuario cuando sea necesario.

\subsection{Resumen de la sección}

Esta clase trata dos mecanismos complementarios de planeación en Codex: uno es el Plan Mode explícito, y el otro es la herramienta \texttt{update\_plan} durante la ejecución. El primero se orienta a la colaboración y a aclarar los requisitos; el segundo, a la gestión del progreso en tiempo de ejecución.
\section{Plan Mode: generación colaborativa de especificaciones}

La primera forma de planeación de Codex es Plan Mode. El usuario puede entrar a este modo con \texttt{/plan}. Al entrar, el sistema inserta en el developer message un conjunto completo de restricciones de comportamiento, de modo que el modelo primero explore, luego aclare y después forme un plan en Markdown que se pueda ejecutar.

\begin{figure}[H]
\centering
\includegraphics[width=\textwidth]{frames/frame-00-03-35.jpg}
\caption{El núcleo de Plan Mode no es ejecutar de inmediato, sino formar una especificación realizable mediante la conversación.\protect\footnotemark}
\end{figure}
\footnotetext{Intervalo de tiempo en el video: 00:03:25--00:03:45.}

Plan Mode se combina con una herramienta especial: \texttt{request\_user\_input}. Esta herramienta solo está disponible en Plan Mode y sirve para plantearle al usuario de una a tres preguntas breves, esperar su respuesta y así reducir los errores de implementación causados por requisitos poco claros.

\begin{knowledgebox}{¿Dónde queda el producto de Plan Mode?}
El docente enfatiza en particular que lo que se forma al terminar Plan Mode es un texto en Markdown, pero que no se guarda automáticamente como un archivo local. Existe dentro de la context window, como contexto para la etapa de ejecución posterior. Dicho de otro modo, es una especificación dentro del contexto, no un documento persistente en el repositorio.
\end{knowledgebox}

Las restricciones de Plan Mode también son estrictas: en esta etapa se debe hacer una exploración que no modifique nada, sin editar el código directamente. Se parece más a una sesión de aclaración de requisitos y diseño de especificaciones, al final de la cual se le pregunta al usuario si desea ejecutar el plan. Una vez que el usuario confirma, el sistema regresa al modo de ejecución predeterminado.

\subsection{Escenarios de uso de Plan Mode}

Plan Mode es adecuado para tareas cuyos límites de requisitos aún no están claros, cuya ruta de implementación tiene varias bifurcaciones, o que requieren que el usuario elija entre distintas alternativas. Por ejemplo, una refactorización grande puede involucrar compatibilidad de API, migración de datos y estrategia de pruebas; en ese caso, dejar primero que el modelo converja la especificación mediante preguntas es más estable que ponerse a trabajar directamente.

\begin{warningbox}{Plan Mode no es una herramienta persistente de gestión de proyectos}
El plan de Plan Mode es solo un objeto de contexto. Puede guiar la sesión actual, pero si se desea guardarlo, recuperarlo o auditarlo entre sesiones, de todas formas se debe usar un plan en archivo, como en este repositorio, por ejemplo \texttt{task\_plan.md}, \texttt{NOTE\_GENERATION\_TODO.md}, etcétera.
\end{warningbox}

\subsection{Resumen de la sección}

Plan Mode es un modo de colaboración. Mediante la developer instruction y \texttt{request\_user\_input} restringe el comportamiento del modelo, y convierte en un flujo explícito el principio de convertir en producto la idea de «preguntar bien antes de escribir código».

\section{update\_plan: el modelo de progreso en vivo durante la ejecución}

La segunda forma de planeación es \texttt{update\_plan}. No es un interruptor para entrar o salir del Plan Mode, sino una herramienta de tiempo de ejecución que mantiene la checklist o lista de pendientes de la tarea actual. El expositor la llama un live progress model, es decir, la representación interna que el agent tiene del estado de avance actual, la cual también puede mostrarse al usuario.

\begin{figure}[H]
\centering
\includegraphics[width=\textwidth]{frames/frame-00-06-30.jpg}
\caption{El schema de \texttt{update\_plan} es muy pequeño: un campo de explicación más una lista de pasos con estado.\protect\footnotemark}
\end{figure}
\footnotetext{Intervalo de tiempo en el video: 00:06:20--00:06:40.}

El schema de \texttt{update\_plan} es muy sencillo:

\begin{lstlisting}[caption={Estructura conceptual de update\_plan}]
{
  "explanation": "Why the plan is being updated",
  "plan": [
    {"step": "Inspect source files", "status": "in_progress"},
    {"step": "Patch implementation", "status": "pending"},
    {"step": "Run verification", "status": "pending"}
  ]
}
\end{lstlisting}

Cada step tiene tres estados posibles: \texttt{pending}, \texttt{in\_progress}, \texttt{completed}. El sistema impone la restricción de que a lo sumo un step puede estar en \texttt{in\_progress}, lo que garantiza que el plan siempre tenga un foco actual claro.

\begin{importantbox}{Una misma herramienta sostiene tres dinámicas}
\texttt{update\_plan} sostiene a la vez la creación del plan, la actualización de estado y el replan. No existen API independientes como \texttt{createPlan}, \texttt{modifyStep} o \texttt{replan}; en cada llamada se envía una instantánea completa y nueva del plan.
\end{importantbox}

A esto se refiere el expositor cuando insiste en las ``plan dynamics'': la primera llamada equivale a crear el plan; las llamadas posteriores pueden hacer avanzar un step de \texttt{pending} a \texttt{in\_progress} o a \texttt{completed}; y si durante la ejecución se descubre que el plan original no es viable, también puede enviarse una nueva instantánea del plan, explicando el motivo en \texttt{explanation}.

\subsection{Resumen de la sección}

Lo esencial de \texttt{update\_plan} no es una API compleja, sino un modelo de estado minimalista. El modelo, al enviar instantáneas del plan de manera repetida, completa la creación, la actualización y el replan dentro del mismo agent loop.

\section{Por qué este diseño es elegante}

El expositor llama a este diseño de Codex un sistema de herramientas general, unificado y mínimo. La razón es que una interfaz de función muy pequeña puede expresar un comportamiento de planificación complejo, mientras que la complejidad queda a cargo de la capacidad semántica del modelo.

\begin{figure}[H]
\centering
\includegraphics[width=\textwidth]{frames/frame-00-09-35.jpg}
\caption{El expositor compara el schema de herramientas de Codex y enfatiza su unificación y su minimalismo.\protect\footnotemark}
\end{figure}
\footnotetext{Intervalo de tiempo en el video: 00:09:25--00:09:45.}

Este diseño tiene tres ventajas:

\begin{enumerate}
    \item \textbf{Interfaz pequeña}: la capa de herramientas solo expone un \texttt{update\_plan}, sin necesidad de agregar una API para cada tipo de operación sobre el plan.
    \item \textbf{Semántica fuerte}: crear, actualizar y replanificar se distinguen por la semántica de los parámetros que interpreta el modelo.
    \item \textbf{Visible para el usuario}: el plan no queda oculto dentro del modelo, sino que puede mostrarse como retroalimentación de progreso.
\end{enumerate}

\begin{knowledgebox}{Instantánea, no parche}
\texttt{update\_plan} se parece más a enviar una «instantánea del estado actual del plan» que a mandar un parche sobre un step determinado. Esto permite que el modelo reordene los pasos en conjunto al replanificar, y también facilita que la interfaz de usuario o el sistema de registro rendericen el estado actual.
\end{knowledgebox}

\subsection{Relación con los planes en archivo}

Cabe notar que \texttt{update\_plan} sigue siendo una herramienta de progreso a nivel de contexto. Para tareas largas, tareas que cruzan varias sesiones y tareas de generación por lotes, los planes en archivo siguen siendo necesarios. El plan dentro del contexto es adecuado para la ejecución inmediata; el todo en archivo es adecuado para la recuperación, la auditoría y la colaboración en equipo.

\subsection{Resumen de la sección}

El diseño de Codex comprime la capacidad de planificación en una herramienta mínima pero suficientemente expresiva. A medida que aumenta la capacidad de modelos como GPT-5.4, el modelo puede usar esta herramienta con mucha fluidez para mantener el estado del plan.
\section{Recomendaciones prácticas: cuándo usar cada tipo de plan}

Combinando el contenido de esta clase, se puede repartir la función de los dos mecanismos así:

\begin{table}[H]
\centering
\begin{tabular}{lll}
\toprule
Mecanismo & Uso principal & Escenario típico \\
\midrule
Plan Mode & Aclarar requisitos, formar la especificación & El objetivo del usuario es ambiguo, hay varias rutas posibles para la tarea, se necesita elegir entre alternativas \\
\texttt{update\_plan} & Gestión del avance durante la ejecución & Ya se comenzó la tarea, se necesita mostrar el estado y ajustarlo dinámicamente \\
todo en archivo & Persistencia entre sesiones & Tareas largas, tareas por lotes, se necesita recuperación y auditoría \\
\bottomrule
\end{tabular}
\caption{Reparto de funciones entre los tres soportes de planificación.}
\end{table}

Para un coding agent, el flujo de trabajo más estable suele ser: primero explorar la base de código, y si es necesario entrar a Plan Mode para aclarar la especificación; durante la ejecución usar \texttt{update\_plan} para gestionar el avance actual; y si la tarea abarca un periodo muy largo, mantener además un plan en archivo.

\begin{warningbox}{Un plan no es mejor entre más detallado}
El valor de un plan está en reducir la incertidumbre, no en generar una sensación de formalidad. Un plan demasiado detallado hace que el agent desperdicie contexto en pasos de bajo valor; un plan demasiado general, en cambio, no logra guiar la ejecución. Un buen plan debe capturar únicamente los pasos que en verdad afectan el orden de ejecución y el control de riesgos.
\end{warningbox}

\subsection{Resumen de la sección}

Plan Mode, \texttt{update\_plan} y el todo en archivo son tres niveles. Cada uno resuelve, respectivamente, el problema de aclarar la colaboración, el del estado de ejecución y el de la persistencia a largo plazo.

\section{Síntesis y ampliación}

La conclusión central de esta clase puede resumirse en una frase: la capacidad de planning de Codex no es un botón único, sino un conjunto de mecanismos por capas. Plan Mode se encarga de aclarar la especificación antes de actuar; \texttt{update\_plan} se encarga de mantener el live progress durante la ejecución; y el todo en archivo se encarga de desacoplar el estado de largo plazo de la context window.

Al final, quien dio la clase señala que, a medida que la capacidad de los modelos aumenta, una herramienta tan minimalista como \texttt{update\_plan} ya basta para que el modelo lleve a cabo creación, avance de estado y replan complejos. Lo que en verdad importa no es el número de herramientas, sino si el schema de la herramienta está alineado con la capacidad semántica del modelo.

\begin{importantbox}{Takeaway final}
Un buen agent harness no necesariamente tiene que acumular muchas herramientas especializadas. Una herramienta lo bastante general, con semántica clara y estado visible, combinada con un modelo fuerte, suele cubrir un espacio de conducta mucho mayor. El \texttt{update\_plan} de Codex es el representante de este tipo de diseño.
\end{importantbox}

Para mantener nuestro repositorio de notas de curso también podemos tomar esta idea como referencia: para tareas cortas, avanzar con un plan en el contexto; para tareas largas, fijar el estado con un todo en archivo; y cada vez que nos encontremos con un bloqueo de descarga, subtítulos, transcripción, compilación, etc., no quedarnos solo en la memoria oral, sino escribirlo de vuelta en la cola, para formar un flujo de trabajo recuperable.

\end{document}
